\subsection{Container shape and dimension}\label{sec:dimension}

Theorems~\ref{thm:selection} and~\ref{thm:oblivious}, and now
Proposition~\ref{prop:n3}, hold for arbitrary compact containers in
$\R^d$. For spherical containers the negative results also transfer
to higher dimensions, for a different reason: each relevant sibling
query involves at most three balls.

\begin{lemma}[Dimension reduction for sibling queries]\label{lem:dimreduce}
Let $k\ge2$ and $d\ge k-1$. Balls of positive radii $a_1,\ldots,a_k$
pack in a ball of radius $R$ in $\R^d$ if and only if they pack in
one of the same radius in $\R^{k-1}$. A single ball fits if and only
if its radius is at most $R$, in every positive dimension.
\end{lemma}
\begin{proof}
The forward embedding from lower to higher dimension preserves all
center distances and norms. Conversely, center the container at zero
and let $c_1,\ldots,c_k$ be the centers of a packing. Their affine hull
$A$ has dimension at most $k-1$. Let $q$ be the orthogonal projection
of zero onto $A$. Since $q$ is orthogonal to $A-q$,
\[
 \|(c_i-q)-(c_j-q)\|=\|c_i-c_j\|\ge a_i+a_j,
 \qquad \|c_i-q\|^2=\|c_i\|^2-\|q\|^2\le(R-a_i)^2.
\]
Here $R-a_i\ge0$ by the original containment. A common translation
of the centers therefore preserves separation and does not increase
their norms. Identify $A-q$ isometrically with a subspace of
$\R^{k-1}$. This translates the whole affine hull; it does not project
the individual centers onto a predetermined lower-dimensional space.
\end{proof}

\begin{corollary}[Transfer of the sharpness results]\label{cor:dimsharp}
For spherical pans in every dimension $d\ge2$, the following hold with
the same radii and widths as in the plane:
\begin{enumerate}
\item placement obliviousness holds for at most three pieces and fails
  at four, as in Theorem~\ref{thm:n4};
\item the twin instances of Theorem~\ref{thm:twins} exclude every
  deterministic state-based rule, and every randomized state-based
  rule fails on some instance with probability at least $1/2$;
\item the rigid family of Theorem~\ref{thm:rigidfloor} has infimum
  exactly $T$, not attained at positive width;
\item the dimension-dependent geometric threshold satisfies
  $1\le\tau_d\le\varphi$, and its restriction to inventories of
  at most four rings satisfies $\tau_{4,d}=\varphi$, unattained.
\end{enumerate}
\end{corollary}
\begin{proof}
By Lemma~\ref{lem:dimreduce}, every query of at most three siblings
in a ball has the same answer in all dimensions $d\ge2$. A fixed
assignment forest whose containers have at most three children can
be realized by assembling these local packings from the leaves:
each child's entire subtree stays inside its outer ball.

The positive part of (1) is Proposition~\ref{prop:n3}. Its negative
part uses the root pair $\{10,5\}$, the hole pair $\{4.9,4.8\}$,
and the forbidden root triple $\{10,4.9,4.8\}$. All are preserved.
The four executions and the two witnesses of Theorem~\ref{thm:twins}
also use only individual, pair and triple queries. The state at the
arrival of the $5$ is identical in the two inventories, and the
necessary choices are still opposite. If the probability of choosing
the root is $p$, the failure probabilities are $1-p$ and $p$.

In (3), conditions (F1) and (F2) are scalar, and (F3) is precisely
a triple query. Thus the entire admissible parameter family and its
function $\rho$ are unchanged, including the approximating sequence.
Finally, the witness of Theorem~\ref{thm:golden} has the root triple
$\{\varphi,s_1,s_2\}$ and the nested pivot $1$. Its failing execution
only queries pairs and triples as well. It gives the same failures
with $\rho=\varphi+3\varepsilon\to\varphi$. The lower bound in (4)
is Theorem~\ref{thm:oblivious} under weak superincreasingness.
For the four-ring lower bound, the forest reduction in
Theorem~\ref{thm:fourfloor} is dimension-independent. Each of its
two remaining cases uses only pairs and triples in spherical
containers, so the rigid floor and the diametral-pocket argument
apply unchanged. Every four-ring failure therefore still has
$\rho>\varphi$; the golden family gives the matching infimum.
\end{proof}

The corollary does not imply $\tau_d=\tau_2$ or $\tau_d=\varphi$:
other counterexamples may involve four or more siblings. Nor does
monotonicity of packing feasibility in the dimension, by itself,
give monotonicity of the threshold for failure of a greedy execution.

\begin{proposition}[The first dimensional separation]\label{prop:dimfour}
There are four distinct sibling radii that fit in the unit ball
in $\R^3$ but not in the unit disk in $\R^2$, namely
$\{441,440,439,438\}/1000$.
\end{proposition}
\begin{proof}
Set $u=8/25$ and use the tetrahedral centers
\[
(u,u,u),\quad(u,-u,-u),\quad(-u,u,-u),\quad(-u,-u,u).
\]
Each norm squared is
$3u^2$ and each pair distance squared is $8u^2$. Even with all
radii enlarged to $441/1000$, containment and separation follow from
\[
 (1-441/1000)^2-3u^2=5281/10^6>0,\qquad
 8u^2-(882/1000)^2=10319/250000>0.
\]
For planar infeasibility, shrink all radii to $r=438/1000$.
The centers would lie in radius $1-r$. A center at zero is impossible
since $1-r<2r$. Otherwise two consecutive centers in angular order
have angular separation at most $\pi/2$, so their squared distance
is at most $2(1-r)^2<(2r)^2$; the difference is $16961/125000$.
This contradicts separation. This is a difference of sibling
feasibility, not a new failure of placement obliviousness.
\end{proof}

\subsection{Square pans: confinement and twin instances}\label{sec:square}

Only the root changes to $[0,s]^2$; every ring's hole is still a disk.
The capacity used by best fit at the root is its inradius $s/2$.
The following sufficient criterion excludes all placements of a trio,
without assuming an optimal tangency pattern.

\begin{lemma}[Quadrant confinement]\label{lem:squareQ}
Let $a\ge b\ge c>0$, $a\le s/2$, and $p\ge q\ge0$. If
\begin{align}
 p&\ge s/2-b, & p^2+(s-a-b)^2&<(a+b)^2,\label{eq:squareQb}\\
 q^2+(s-a-c)^2&<(a+c)^2, &p+q&\ge s-b-c,\label{eq:squareQc}\\
 L:=s-c-a-q&\ge0, &2L^2&<(b+c)^2,\label{eq:squareQbox}
\end{align}
then disks of radii $a,b,c$ do not pack in the square of side $s$.
\end{lemma}
\begin{proof}
Reflect the whole packing about either midline as needed, so the
largest center $A$ has both coordinates in $[a,s/2]$. The coordinates
of $B,C$ lie in $[b,s-b]$ and $[c,s-c]$, respectively. Each absolute
coordinate difference between $A$ and $B$ is at most $s-a-b$.
If $|B_x-A_x|\le p$, the first quadratic inequality contradicts
separation. The negative alternative $B_x-A_x<-p$ is excluded by
$A_x-B_x\le s/2-b\le p$. Thus $B_x>A_x+p$, and likewise
$B_y>A_y+p$, whence $A_x,A_y<s-b-p$.

The same argument with the other quadratic inequality gives
$|C_x-A_x|>q$. Its negative alternative is impossible because
$A_x-C_x<s-b-p-c\le q$. Therefore $C_x>A_x+q$ and
$C_y>A_y+q$. Since $p\ge q$ and $b\ge c$, both small centers lie
in the box $[a+q,s-c]^2$. Their squared distance is at most $2L^2$,
contradicting~\eqref{eq:squareQbox}. All separations used for
packings are non-strict; tangencies have not been excluded.
\end{proof}

\begin{proposition}[An exact square counterexample]\label{prop:squareD}
With side $s=4.8568$, width $w=0.16$ and radii
$\{a,1,b,c\}=\{1.845,1,0.844,0.841\}$, all four rings fit but
best fit places exactly three. The dominant tail is $\rho=337/200$.
Every decimal here and below denotes the exact terminating rational.
\end{proposition}
\begin{proof}
Apply Lemma~\ref{lem:squareQ} with $p=1.591$, $q=1.581$.
The margins in its five inequalities are, in their displayed order,
\[
 \frac{33}{5000},\quad\frac{2079}{25000000},\quad
 \frac{66559}{25000000},\quad\frac1{5000},\quad
 \frac{53587423}{25000000},\qquad L=\frac{2949}{5000}.
\]
The witness places $a$ at $(a,a)$ and $1$ at $(s-1,s-1)$:
$2(s-a-1)^2-(a+1)^2=16337/25000000>0$. In the hole of $a$,
of radius $H=a-w=b+c=1.685$, place the disk of radius $b$ at
relative center $(-c,0)$ and that of radius $c$ at $(b,0)$.

Best fit instead nests $1$ in $a$, since $1\le H<s/2$.
The disk $b$ then fits only at the root: $1+b>H$ and $b>1-w$.
Root feasibility follows by shrinking the disk $1$ in the witness.
At $c$, the root triple is excluded by the lemma, and all holes are
excluded by $1+c>H$, $c>1-w$ and $c>b-w$. The tails are
$179/123$, $337/200$, and $841/844$, so the middle one dominates.
Together with Proposition~\ref{prop:n3}, this establishes the same
three/four-piece transition for a square.
\end{proof}

\begin{theorem}[Square twins]\label{thm:squaretwins}
Let $s=5.1214$, $w=0.301$, and
\[
 I_1=\{2,1,0.850,0.849\},\qquad
 I_2=\{2,1,0.950,0.750\}.
\]
Both full inventories are feasible. In $I_1$ the pivot $1$ must be
placed at the root for the greedy to succeed; in $I_2$ it must be
nested in $2$. Hence no state-based rule succeeds universally, and
every randomized state-based rule has failure probability at least
$1/2$ on one of these instances.
\end{theorem}
\begin{proof}
The hole of $2$ has radius $H=1.699$. Lemma~\ref{lem:squareQ}
excludes $\{2,0.850,0.849\}$ using $(p,q)=(1.72,1.718)$, and
excludes $\{2,1,0.750\}$ using $(p,q)=(2.12,1.39)$; direct rational
substitution verifies all its inequalities. Shrinking the last disk
therefore excludes $\{2,1,z\}$ for every $z\ge0.750$.
On the other hand $\{2,0.950,0.750\}$ has the explicit centers
\[
 (2.96,2),\qquad(0.95,4.1714),\qquad(4.3714,4.3714).
\]
Containment is immediate. The three positive separation margins are
\[
\frac{1314449}{25000000},\qquad\frac{663599}{12500000},\qquad
\frac{221399449}{25000000}.
\]

The root pair $\{2,1\}$ fits in opposite corners since
$2(s-3)^2>9$. Shrinking its unit disk at the same center also
packs $\{2,z\}$ for every $0<z<1$, justifying admission of the
first small ring when the pivot is nested.
This pair, with the two small rings inside $2$,
witnesses $I_1$, as $0.850+0.849=H$.
For $I_2$, use the displayed root triple and nest $1$ in $2$.
Every small ring fits alone in $H$, but none fits there beside $1$,
since even $1+0.750>H$. None fits in the pivot's hole of radius
$0.699$, and the last cannot nest in the preceding small ring:
$0.849>0.850-w$ and $0.750>0.950-w$.

If the pivot is nested, both small rings are forced to the root,
where the first triple is infeasible and the second feasible.
If the pivot is at the root, both are excluded from the root and
must use $H$, where only the first pair fits
($0.950+0.750=1.700>H$). All subsequent decisions are therefore
forced. At the pivot, the state, side, width and inventory size are
identical. If the root is chosen with probability $p$, the failure
probabilities are $1-p$ and $p$. The dominant tails are $1.699$ and
$1.700$, respectively.
\end{proof}

\subsection{An algebraic upper bound for the square threshold}
\label{sec:squarelimit}

Define $\tau_{\square}$ as the infimum of $\rho$ over square-pan
instances admitting a failing greedy execution, with side and positive
width allowed to vary, as in the definition of $\tau$.

\begin{theorem}[An approximating square family]\label{thm:squarelimit}
Let $Y$ be the unique positive root of
\[
 (17+10\sqrt2)Y^2+(72+16\sqrt2)Y-(112+96\sqrt2)=0.
\]
Then
\[
 1\le\tau_{\square}\le Y=1.684487745872346\ldots.
\]
There are four-ring counterexamples to best fit with strictly
decreasing rational radii, positive rational width and rational side,
whose values of $\rho$ tend to $Y$ from above. No optimality of $Y$
is asserted, either globally or for all parameters of the confinement
criterion.
\end{theorem}
\begin{proof}
Put $k=\sqrt2$, $C=1+k/2$, and consider first the balanced data
$a_0=1+t$, $b_0=c_0=t$, $s_0=C(2+t)$,
$p_0=q_0=s_0/2-t$. The common quadratic margin is
\[
 G(t)=(a_0+t)^2-p_0^2-(s_0-a_0-t)^2
 =\frac{(17+10k)t^2+(36+8k)t-(28+24k)}8.
\]
It is strictly increasing for $t\ge0$ and has exactly one positive
root $t_0=Y/2$. Indeed,
\[
 G(21/25)=\frac{8897}{5000}-\frac{639k}{500}<0,
 \qquad
 G(17/20)=\frac{5953}{3200}-\frac{399k}{320}>0,
\]
using $7/5<k<10/7$. For $21/25<t<17/20$,
\[
 L_0=s_0/2-a_0=k/2+(k-2)t/4,\qquad
 \tfrac12<\tfrac{229}{400}<L_0<\tfrac57,
 \qquad 2L_0^2<\tfrac{50}{49}<4t^2.
\]
Thus containment of $a_0$, positivity of $p_0$, and the final
confinement inequality have strict margins.

Fix any $t_0<t<17/20$ and perturb by $\eta>0$:
\begin{gather*}
 a=1+t+2\eta,\quad m=1,\quad b=t+\eta,\quad c=t-\eta,\\
 w=1-t+2\eta,\quad H=a-w=2t,\quad
 s=C(a+1),\quad p=q=s/2-t.
\end{gather*}
At $\eta=0$ both quadratic margins equal $G(t)>0$. The strict
inequalities just established, including $a<s/2$ and $p>0$,
are continuous in $\eta$; finitely many such inequalities remain
strict on some interval $0<\eta<\eta_0(t)$. The two remaining
walls hold identically:
\[
 p-(s/2-b)=\eta>0,\qquad p+q-(s-b-c)=0.
\]
Choose also $\eta<t$, $\eta<1-t$, and $3\eta<2t-1$.
Then $a>1>b>c>w>0$, and Lemma~\ref{lem:squareQ} excludes the
root triple $\{a,b,c\}$.

The witness places $\{a,1\}$ in opposite corners and $\{b,c\}$
on a diameter of the large hole. Both are legal because
$2(s-a-1)^2=(a+1)^2$, both root radii are at most $s/2$, and
$b+c=H$. Best fit nests $1$, since $1<H<a<s/2$.
The differences excluding the alternative holes for the two
small rings are
\begin{gather*}
 1+b-H=1-t+\eta>0,\quad 1+c-H=1-t-\eta>0,\\
 b-(1-w)=3\eta>0,\quad c-(1-w)=\eta>0,\quad
 c-(b-w)=1-t>0.
\end{gather*}
The first small ring fits at the root by shrinking the pivot in
the witness; the last is excluded everywhere. Best fit places
exactly three rings. Moreover
\[
 \frac{1+b+c}{a}=\frac{1+2t}{1+t+2\eta}<2t,
 \qquad \frac cb<1<2t,
\]
where the first inequality uses $1<2t^2+4t\eta$ and $t>21/25$.
Thus $\rho=2t\to2t_0=Y$ from above, proving the upper bound.
The lower bound follows from Theorem~\ref{thm:oblivious}.

For entirely rational instances choose rational $t\downarrow t_0$
and then rational $\eta$ in the nonempty interval above. Replace
$s=C(a+1)$ by a slightly larger rational and reset $p=q=s/2-t$.
All strict inequalities persist by continuity in $s$, the two
displayed wall identities persist exactly, and the root pair gains
strict separation. Density of the rationals completes the construction.
The unperturbed balanced limit itself is not a strict-radius
counterexample and is not used as one.
\end{proof}

Eliminating the radical shows that $Y$ also satisfies
\[
 89Y^4+1808Y^3+4704Y^2-9984Y-5888=0.
\]
The positive branch is specified by the quadratic over $\mathbb Q(\sqrt2)$,
not by an arbitrary choice of a positive root of this quartic.
Exact rational examples in the accompanying code attain
$\rho=1.684488$, $1.684487746$, and $1.6844877458724$.
These finite checks certify only those instances; the limiting bound
uses the continuity proof above.

\paragraph{Formal verification scope.}
\texttt{Square.lean} proves Lemma~\ref{lem:squareQ} as a universal
Cartesian exclusion over linearly ordered commutative rings, including
the reflections. \texttt{SquareTwins.lean} proves the required twin
exclusions and the explicit feasible triple with the same Cartesian
predicate. \texttt{SquareLimit.lean} checks the balanced-gap and norm
polynomial identities. These files contain 24 theorems, in addition
to the 57 earlier algebraic certificates and the three in
\texttt{FourRing.lean}. The Euclidean interpretation,
assignment forests, algorithmic arguments, dimension reduction and
continuity remain written proofs; they are not claimed as fully
formalized Lean theorems. All new numerical checks use exact rational
arithmetic, and all displayed decimals specifying instances are exact.
